% =============================================================================
%  Migrating GEOUNED to pyOCC — Technical report
%
%  Self-contained document. Build with:
%      pdflatex pyocc_migration_report_en.tex
%      pdflatex pyocc_migration_report_en.tex   (second pass: TOC + refs)
%
%  Synthesised from the project working log CLAUDE.md.
%  Branch: georeverse-migration. Date: 2026-09-06.
%  English translation of docs/reports/pyocc_migration_report.tex.
% =============================================================================
\documentclass[11pt,a4paper]{article}

\usepackage[utf8]{inputenc}
\usepackage[T1]{fontenc}
\usepackage{lmodern}
\usepackage[english]{babel}
\usepackage[a4paper,margin=2.4cm]{geometry}
\usepackage{textcomp}
\usepackage{xcolor}
\usepackage{booktabs}
\usepackage{longtable}
\usepackage{array}
\usepackage{ragged2e}
\usepackage{enumitem}
\usepackage{parskip}
\usepackage{microtype}
\usepackage[hidelinks]{hyperref}
\setlength{\emergencystretch}{3em}

% --- style -------------------------------------------------------------------
\definecolor{accent}{HTML}{1F6F8B}
\hypersetup{
  pdftitle={Migrating GEOUNED to pyOCC - Technical report},
  pdfsubject={GEOUNED CAD backend refactor},
  colorlinks=true, linkcolor=accent, urlcolor=accent, citecolor=accent
}

% Literal identifiers/paths: \detokenize prints _, #, &, %, ~ verbatim.
\providecommand{\code}[1]{\texttt{\detokenize{#1}}}
\newcommand{\statusok}{\textbf{\textcolor{accent}{repaired}}}
\newcommand{\statusdone}{\textbf{\textcolor{accent}{resolved}}}
\newcommand{\statusnone}{\textit{no repair}}
\newcommand{\statuslimit}{\textit{accepted limitation}}

\newcolumntype{L}[1]{>{\RaggedRight\arraybackslash}p{#1}}

\setlist[itemize]{leftmargin=1.4em, itemsep=2pt, topsep=3pt}
\setlist[enumerate]{leftmargin=1.6em, itemsep=4pt, topsep=3pt}

\title{Migrating GEOUNED to pyOCC\\[6pt]
       \large\normalfont Technical report on the CAD backend refactor}
\author{GEOUNED project \textemdash{} branch \texttt{georeverse-migration}}
\date{September 6, 2026}

% =============================================================================
\begin{document}
\maketitle

\begin{center}
\small
\begin{tabular}{@{}ll@{}}
\toprule
Project             & GEOUNED-org/GEOUNED \\
Branch              & \code{georeverse-migration} \\
Default engine      & \code{ocp} (OCP / pybind11) \\
Supported engines   & \code{freecad} $\cdot$ \code{occ} $\cdot$ \code{ocp} \\
Scope               & CadToCsg (GEOUNED) + CsgToCad (GEOReverse) \\
Source              & working log \code{CLAUDE.md} \\
\bottomrule
\end{tabular}
\end{center}

\vspace{4pt}
\noindent\fbox{\parbox{\dimexpr\textwidth-2\fboxsep-2\fboxrule}{%
\small\centering
\textbf{3\,/\,3} CAD engines with the test suite green \quad$\vert$\quad
\textbf{0} lost particles (\code{test_models} corpus, 143 d1suned runs) \quad$\vert$\quad
\textbf{93.6\%} tallies within 2$\sigma$ / \textbf{0\%} above 3$\sigma$ \quad$\vert$\quad
\textbf{0} direct \code{Part}/\code{FreeCAD} imports in GEOUNED%
}}

\vspace{6pt}
\tableofcontents
\vspace{6pt}
\hrule

% =============================================================================
\section{Executive summary}

GEOUNED converts CAD geometry into CSG geometry (and back) for the Monte Carlo
transport codes \textbf{MCNP, OpenMC, PHITS and Serpent}. It historically
depended entirely on FreeCAD's \code{Part} API and, in particular, on
\code{BOPTools.SplitAPI.slice} for solid decomposition. That coupling has been
removed.

The work done falls into five blocks:

\begin{itemize}
  \item \textbf{Kernel isolation.} All code that touches FreeCAD/OCCT now lives
        in a single package, \code{geo}, behind a stable name-level interface.
        No other file in GEOUNED imports \code{Part} or \code{FreeCAD}.
  \item \textbf{Three geometry engines.} \code{geo} has three interchangeable
        implementations selected via the environment variable
        \code{GEOUNED_CAD_ENGINE}: FreeCAD (\code{freecad}), pythonocc-core /
        SWIG (\code{occ}) and OCP / pybind11 (\code{ocp}, currently the default).
  \item \textbf{Quantitative verification.} Every fix has been validated by two
        independent methods: point-by-point checking of \code{check_sign}
        against independently built CAD geometry, and a stochastic volume check
        that runs the resulting MCNP model in d1suned (tally $\approx$ 1.0).
  \item \textbf{Data-driven geometry debugging.} Dozens of real bugs
        \textemdash{} many predating the migration \textemdash{} found and fixed
        by following concrete traces: AND/OR signs in composite surfaces,
        tangency degeneracies in \code{Gsplit}, surface classification,
        tolerance scaling and residual \emph{slivers}.
  \item \textbf{CAD defect repair.} A new load-time layer that detects corrupt
        solids (pathologically short edges, duplicated boundary rings, phantom
        cuts) and repairs them automatically when that is fast and reliable, or
        flags them untranslated when it is not.
\end{itemize}

The reverse pipeline, \textbf{GEOReverse} (CsgToCad), which was initially out of
scope because of its hard dependency on the \code{.FCStd} format, has received
its own migration on the current branch: it is engine-switchable, exports real
STEP via XCAF under the pyOCC engines, and colours each solid by material.

Status as of this report: all three engines pass their suites; the reference
corpus \code{Solidos/test_models} translates with no lost particles and no
failures above 3$\sigma$; \code{occ} and \code{ocp} produce identical results
over the whole corpus. Several GEOReverse items and a few exotic surfaces
remain unported (Section~\ref{sec:pending}).

% =============================================================================
\section{Context and motivation}

\subsection{The problem that triggered the migration}

FreeCAD's boolean \code{Cut} can \textbf{silently} fail to split a solid into
two when the cutting plane's intersection coincides exactly with a pre-existing
tangency line (for example, a plane cutting the corner where a cylinder is
already tangent to two cube faces). Symptom: \code{Cut} returns the original,
uncut solid instead of raising an error. Root cause: the resulting fragments
only touch along a line (zero-area contact), producing a non-manifold shape that
OCC's BOP algorithm does not cleanly resolve. Round-tripping through a STEP
export/import ``fixes'' it because the STEP healer re-splits the compound.

The proper fix \textemdash{} building a face-adjacency graph that excludes edges
shared by more than two faces and reconstructing one solid per connected
component \textemdash{} requires fine-grained topological control that FreeCAD's
\code{Part} API, itself a wrapper over OCCT, does not expose conveniently.

\subsection{Decision: pythonocc-core, not CadQuery}

\begin{itemize}
  \item Fine-grained topological control is needed (face/edge adjacency graphs,
        \code{BRepCheck_Analyzer}, \code{ShapeFix_Shape}, tolerance control over
        \code{BRepAlgoAPI_*} / \code{BOPAlgo_Splitter}) for exactly the
        degenerate cases above.
  \item Migrating to pyOCC \emph{removes} a layer (the \code{Part} wrapper)
        rather than replacing it with a different one (CadQuery).
  \item CadQuery is a better fit for high-level parametric construction, not the
        low-level topology surgery GEOUNED's decomposition requires.
\end{itemize}

% =============================================================================
\section{Architecture: the \texorpdfstring{\code{geo}}{geo} package}

\subsection{Attempt 1 (retired): the \texorpdfstring{\code{GeometryBackend}}{GeometryBackend} ABC}

The first design was a classic Adapter / Ports \& Adapters pattern: an abstract
base class \code{GeometryBackend} with neutral dataclasses (\code{GSolid} /
\code{GFace} / \code{GEdge} / \code{GVector} / \dots) that each carried a
\code{native} object plus a reference to the backend that produced them, and a
concrete \code{FreeCADBackend} implementation. All of GEOUNED called through a
single injected instance (\code{_backend.split(gsolid, tool, tol)}, \dots). It
was fully implemented and validated against the whole suite, but was judged
\textbf{too heavy and indirect} \textemdash{} the dataclasses carrying their own
backend and the \code{_backend.method(x, ...)} convention \textemdash{} for what
it bought. It was deleted entirely along with its tests.

\subsection{Attempt 2 (current): the \texorpdfstring{\code{geo}}{geo} package}

\code{src/geouned/geo/} is the \textbf{only place in GEOUNED allowed to import}
\code{Part} / \code{FreeCAD} / \code{BOPTools}.

\begin{longtable}{@{}L{0.30\textwidth} L{0.64\textwidth}@{}}
\toprule
\textbf{File} & \textbf{Responsibility} \\
\midrule
\endhead
\code{vector_geometry.py} & Pure math with no FreeCAD dependency: \code{GVector},
  \code{GMatrix}, \code{GBoundBox} and converters. \\
\code{surface_geometry.py} & Analytic predicates over surface descriptors
  (\code{is_same_*_surface}, \code{is_inside_*}, \code{find_can_plane},
  plane-plane and line-line intersections). \\
\code{solid_defects.py} & Whole-solid CAD defect detection
  (\code{find_short_edges}, \code{find_split_ring_faces},
  \code{near_surface_pair}). \\
\code{io_utils.py} & Process/loading utilities: \code{GLabelNode} (STEP label
  tree), \code{suppress_native_stdout}. \\
\code{_freecad_impl.py} & FreeCAD implementation: surface/curve classification,
  eagerly built topology classes
  (\code{GEdge}/\code{GWire}/\code{GFace}/\code{GShell}/\code{GSolid}) carrying
  real behaviour as methods, \code{Gmake_*} constructors and
  \code{Gcut}/\code{Gcommon}/\code{Gfuse}/\code{Gsplit} operations. \\
\code{_occ_impl.py} & The same name surface against pythonocc-core (SWIG). \\
\code{_ocp_impl.py} & The same name surface against OCP (pybind11). \\
\code{__init__.py} & Single import point. Reads \code{GEOUNED_CAD_ENGINE} once
  and re-exports the names from the resolved implementation. \\
\bottomrule
\end{longtable}

\paragraph{Interchangeability principle.}
Engine selection happens at import time, not through dependency injection. A new
engine only has to define the same class and function names;
\code{geo/__init__.py} chooses which to re-export and the rest of GEOUNED does
not change. GEOReverse replicates exactly the same pattern with its own
\code{GEOReverse/Modules/__init__.py}.

\subsection{Design of the \texorpdfstring{\code{geo}}{geo} classes}

\begin{itemize}
  \item Faces are eagerly classified into one of 5 analytic surfaces (plane /
        cylinder / cone / sphere / torus) via \code{Gclassify_surface}.
        Composite surfaces (RoundCorner, Can, TCone, MultiPlane,
        MultiRoundCorner, ReversedConeCylinder) are assembled by GEOUNED itself
        with \code{Gcut}/\code{Gcommon}/\code{Gfuse}; \code{geo} never models
        them.
  \item Curve classification is tolerant: it returns \code{None} for a
        degenerate edge or a real-but-unsupported type (e.g.\ hyperbola) rather
        than raising, because solids are built eagerly and an incidental edge
        must not abort the whole solid.
  \item \code{Gsplit()} returns a \code{SplitResult} that can never come back
        empty silently: tangency/degeneracy cases are resolved internally
        (tolerance retry, non-manifold reconstruction, component separation) and
        reported with \code{degenerate_case_handled=True}.
  \item \code{utils/geometry_gu.py::FaceGu} / \code{SolidGu} \textemdash{} a
        second native wrapper predating \code{geo}, used by the decomposition
        code \textemdash{} now \textbf{inherit} from \code{GFace} / \code{GSolid}
        and only layer the decomposition-side extras on top, instead of
        duplicating the construction.
\end{itemize}

% =============================================================================
\section{Decoupling FreeCAD from GEOUNED}

Closing the \code{import Part} / \code{import FreeCAD} statements was done file
by file, tracing every real consumer before touching the source. \code{core.py}
(specifically \code{_set_geometry_bounding_box}) was the last holdout, precisely
because its \code{BoundBox} is consumed outside the migration's original scope
(\code{void.py}, \code{write_files.py}). After it was closed, \textbf{zero
files} in the GEOUNED subpackage import \code{Part}/\code{FreeCAD}/\code{BOPTools}
directly.

\subsection{Unifying duplicated representations}

``A plane'' had three parallel representations: the descriptor \code{geo.GPlane},
the \code{PlaneParams} class in \code{basic_functions_part1.py}, and the
\code{Plane} class in \code{build_region/Objects.py}. They were consolidated:

\begin{itemize}
  \item ``analytic surface + bounding box $\rightarrow$ finite clipped solid''
        for plane/cylinder/cone was unified into three pure functions in
        \code{build_region/Objects.py} (\code{plane_polygon_from_box},
        \code{cylinder_from_box}, \code{cone_from_box}), in
        \code{GVector}/\code{GBoundBox} space, with no duplication.
  \item The \code{Plane}/\code{Cylinder}/\code{Cone}/\code{Sphere} classes in
        \code{Objects.py} collapsed into a single \code{CellSurface} wrapping a
        \code{geo} descriptor; \code{get_surface()} now builds
        \code{GPlane}/\code{GCylinder}/\dots directly via
        \code{.from_values(...)}, with no translation step and no native detour.
  \item The Tier-1 \code{*OnlyParams} classes now store \code{GVector} directly
        instead of \code{FreeCAD.Vector}; \code{check_sign_primitive} collapsed
        to a dict dispatch against the shared \code{is_inside_*} functions.
\end{itemize}

\paragraph{Recorded methodological lesson.}
A static ``who reads this field'' survey is necessary but not sufficient when
the value passes through an intermediate object before reaching the operation
that cares about its type. Several times, a \code{grep}-based survey concluded
that every consumer tolerated \code{GVector}, and the real test suite surfaced
native call sites the \code{grep} missed. Rule: follow every real
\code{AttributeError}/\code{TypeError} from the full suite, not just static
tracing.

\subsection{Two surface-numbering scopes}

Every \code{GeounedSurface} has \code{.bVar} (the signed id it is referred to
by) and, for composite types, \code{.region} (the AND/OR boolean definition) and
\code{.components} (the numbering\,$\leftrightarrow$\,surface relation). They are
assigned in two phases with genuinely different numbering: in
\textbf{decomposition} the \code{.bVar} is a throwaway local id whose only job
is to keep the AND/OR expression consistent for one CAD build; in
\textbf{conversion} it is overwritten with the global, deduplicated canonical
id. The two phases \emph{cannot} read each other's \code{.region}, but they do
share the AND/OR \textbf{rule}, extracted into pure functions
(\code{can_region}, \code{tcone_region}, \code{round_corner_region}, \dots).

% =============================================================================
\section{Writer reorganisation}

The four output-format writers had near-identical implementations of several
methods. They were reorganised along the real card-convention boundary:

\begin{itemize}
  \item \code{write/mcnp_like/} \textemdash{} MCNP, Serpent and PHITS (the three
        MCNP-lineage text formats, sharing full-line/inline comment conventions
        and banners).
  \item \code{write/openmc/} \textemdash{} structurally different (XML/Python
        output, no comment-card concept).
  \item \code{write/functions.py} \textemdash{} shared home for
        \code{get_cell_surf_summary}, \code{simplify_planes},
        \code{sorted_surfaces}, used by all four formats.
\end{itemize}

The \code{CommonInputWriter} mixin (inherited by
\code{McnpInput}/\code{SerpentInput}/\code{PhitsInput}) collects what is
duplicated only among the three MCNP-lineage formats, parametrised by class
attributes (\code{_surface_formatter}, \code{inline_comment_char}, \dots).

\paragraph{A real correction, confirmed with the user.}
The \code{sorted_surfaces} unification fixes an inconsistency: only MCNP applied
\code{Surfaces.IndexOffset} when relabelling each \code{bVar};
Serpent/OpenMC/PHITS skipped it. All four formats now use the MCNP version. The
effect is currently dormant (\code{IndexOffset} is 0 unless
\code{settings.startSurf} is wired, which was also done in this migration) but it
is no longer a divergence between formats.

% =============================================================================
\section{Verification campaigns}
\label{sec:verif}

\subsection{Point-by-point \texorpdfstring{\texttt{check\_sign}}{check\_sign} verification}

For every composite surface found in a model, an \emph{independent} CAD
reference solid is built via the same path \code{Gsplit} itself uses to build
that composite, random points are sampled strictly inside the construction box,
and \code{check_sign(point, surface)} is compared against
\code{gsolid.is_inside(point)}.

\begin{longtable}{@{}L{0.22\textwidth} L{0.15\textwidth} L{0.57\textwidth}@{}}
\toprule
\textbf{Composite type} & \textbf{Result} & \textbf{Notes} \\
\midrule
\endhead
RoundCorner & 300/300 & Many instances across dozens of files. \\
Can (Fwd / Rev) & 300/300 & A sampling artefact (bounding-box padding) was
  identified and fixed: sample inside the construction box, not the resulting
  solid's box. \\
TCone (Fwd / Rev) & 300/300 & Surfaced a real bug in \code{TCone_region} (it
  ignored \code{p1_configuration}/\code{p2_configuration}), fixed. \\
MultiRoundCorner & 300/300 & 17 instances across 13 files. First end-to-end
  validation of this type. \\
MultiPlane & 300/300 & Surfaced that \code{multiplane()} could never detect a
  MultiPlane (instance-vs-class comparison), fixed. \\
ReversedConeCylinder & partial & Structurally unreachable from the decomposition
  scan; verified from the conversion side with one real discrepancy on
  \code{hylife-v06} solid113 (271/300), tied to a known GEOReverse mismatch. \\
\bottomrule
\end{longtable}

\subsection{MCNP stochastic volume check}

A different, more end-to-end question: does a \emph{fully translated} MCNP model,
actually run, reproduce the \textbf{same volumes} as the original CAD solids?
The model is exported with \code{volSDEF=True} (isotropic photon source from an
enclosing sphere + an F4 track-length tally on each cell, normalised by the
CAD-computed volume), run in \code{MODE P} + \code{VOID} in d1suned, and the
expected tally result is exactly \textbf{1.0}. A significant deviation ($>$
3$\sigma$) means the written surface signs do not reconstruct the same solid.

\begin{center}
\small
\begin{tabular}{@{}L{0.68\textwidth} r@{}}
\toprule
\textbf{Corpus \code{Solidos/test_models} (excluding \code{Big_*})} & \textbf{Value} \\
\midrule
Files that translate                                  & 133 / 138 \\
Tallies within 2$\sigma$ of 1.0                        & 93.6\% \\
Marginal tallies (2\textendash 3$\sigma$, expected noise) & 6.4\% \\
Real failures ($>$ 3$\sigma$)                          & 0\% \\
Files with lost particles                              & 0 / 143 runs \\
\code{occ} vs \code{ocp} divergence over the corpus    & byte-for-byte identical \\
\bottomrule
\end{tabular}
\end{center}

The 5 files that do not translate are genuine CAD defects that the detection
layer (Section~\ref{sec:defects}) correctly flags, plus one native-crash case
accepted as a limitation (\code{ConeSphere.stp} under \code{occ}/\code{ocp}).

% =============================================================================
\section{Geometry fixes}

Dozens of real bugs, most predating the migration (confirmed against the last
pre-migration commit), found by following concrete traces \textemdash{} never by
guessing \textemdash{} and verified before and after. Grouped by nature:

\begin{longtable}{@{}L{0.16\textwidth} L{0.56\textwidth} L{0.22\textwidth}@{}}
\toprule
\textbf{Category} & \textbf{Representative examples} & \textbf{Impact} \\
\midrule
\endhead

\textbf{AND/OR signs in composite surfaces} &
\code{can_region} cone + apex-plane Reversed branch used AND where the physical
sign requires OR (two instances); \code{TCone_region} substituted a single
orientation-based value for the real configuration; \code{round_corner_region}
applied \code{fwd_cyl} twice in the \code{p1 == p2} branch;
\code{cone_apex_plane} double-compensated the Reversed case;
\code{add_reversedCC} redesigned
(\code{plane_region AND (s1 AND ... AND sn)}) and its \code{AddPlanes} moved from
unconditional OR to AND; \code{AdjacentMultiplanePlanes} became a list-of-lists
(AND within a group, OR across groups). &
Wrong CSG volume with a valid boolean expression; in several cases ``leaks''
that engulfed half the model. Verified with d1suned after each fix. \\
\addlinespace

\textbf{Degeneracy / tangency in \code{Gsplit}} &
Coaxial cone-cone split fallback (\code{_try_coaxial_cone_split}): analytic
reconstruction of the exact intersection circle, pre-split of that edge and a
retry with fuzzy-tolerance escalation; the implicit-zero \code{splitTolerance}
default moved from 0.0 to $10^{-4}$ under \code{occ}/\code{ocp} (OCCT 7.9.3 is
more sensitive to near-tangency); \code{_repair_non_manifold_solid} now
validates its own reconstruction (validity + volume conservation) before
trusting it. &
Decomposition silently returning 1 piece instead of 2 or 3; fragments with
inflated volume. \\
\addlinespace

\textbf{Surface classification} &
\code{is_same_plane_surface} compared two planes' offsets with anti-parallel
axes as \code{d1 == d2} instead of \code{d1 == -d2}; \code{convex_planes} had a
dead sign test (\code{cross.dot(ref)} always 0) and did not normalise
\code{atan2} to $[0, 2\pi)$; \code{multiplane()} compared an instance against a
class (\code{Gclassify_curve(e) is GLine}), making MultiPlane undetectable;
\code{is_closed_cylinder_cone} accepted a face that was not a closed 360$^\circ$
cylinder $\rightarrow$ a real topological invariant by winding
(\code{_is_closed_by_winding}); \code{cyl_plane_region_conf} used the seed face
for both ends of a split cylinder. &
Faces classified into the wrong composite type; tallies of 0.0 and hundreds of
$\sigma$ deviations, all closed. \\
\addlinespace

\textbf{Tolerances} &
\code{min_area} was never wired into \code{order_plane_face} or the candidate
generators; new \code{min_face_width} + \code{CharacteristicWidth}
(characteristic width of a face via radii of gyration) to detect ``narrow but
large-enough-area'' \emph{slivers}, ported to all three engines;
\code{Tolerances.scaled(volume)} scales \code{min_area}/\code{min_face_width}
down for tiny decomposed fragments; \code{eligible_plane} ignored the
user-configured tolerances. &
Spurious sliver planes entering the cell definition $\rightarrow$ lost
particles. \\
\addlinespace

\textbf{Residual boolean-cut slivers} &
\code{other_face_edge} with a \code{skip_slivers} mode to walk past sliver faces
and find the real one; sliver filter in \code{merge_same_surface_faces} before
the $O(n^2)$ adjacency walk; \code{sort_range} with a ``no-progress'' fallback
for fragments joined only modulo $2\pi$ (was infinite recursion). &
Can/RoundCorner detection failing or hanging on geometry with cut artefacts. \\
\addlinespace

\textbf{Cross-engine} &
The \code{ShapeUpgrade_UnifySameDomain} flag that triggers a native crash is
\emph{opposite} between OCP and pythonocc-core on the same OCCT 7.9.3 geometry;
\code{get_join_cone_cyl} used \code{.Index} where a merged shell has
\code{.Indexes} (\code{freecad} only); MKL's LAPACK conflicted with
\code{occt}/\code{tbb} in the same process $\rightarrow$ switch to OpenBLAS in
\code{pyoccenv} and \code{ocpenv}. &
47 tests crashing on \code{freecad}/\code{occ} with an \code{AttributeError};
native access violations. \\
\addlinespace

\textbf{Non-geometric, found along the way} &
\code{generic_split} double-wrapped an already-wrapped \code{GSolid}
$\rightarrow$ decomposition was silently a no-op; \code{check_intersection} with
a \code{Gdistance} early-exit classified fully nested solids as disjoint
$\rightarrow$ the enclosure void cell did not exclude the spheres inside it
(\code{w_encl.stp}, tally 0.0); \code{expand_regions_to_integer} lacked a
\code{bool} guard that all five cell writers needed and only MCNP had. &
Correctness bugs that degraded decomposition or output invisibly. \\
\bottomrule
\end{longtable}

% =============================================================================
\section{Phased pyOCC migration}

The pyOCC migration was tackled in phases, each with an explicit go/no-go
criterion before writing code at scale.

\begin{enumerate}
  \item \textbf{Phase 1 \textemdash{} Scaffolding and go/no-go test.} Engine
        selection via \code{GEOUNED_CAD_ENGINE}; \code{geo/__init__.py} imports
        from \code{_freecad_impl} or the new \code{_occ_impl}. Only the minimum
        was implemented to test \code{Gsplit} against \code{rev_pipe.stp},
        already characterised as broken in FreeCAD. Result:
        \code{BOPAlgo_Splitter} (OCCT 7.9.3) returns \textbf{3 valid solids} on
        the first attempt with no repair, where FreeCAD returned 1 invalid
        solid. A strong signal to continue.
  \item \textbf{Phase 2 \textemdash{} Full port.} All 5 surface descriptors and
        4 curve descriptors,
        \code{Gclassify_surface}/\code{Gclassify_curve}, full
        \code{GEdge}/\code{GWire}/\code{GFace}/\code{GShell}/\code{GSolid}
        (including \code{find_interior_point}, \code{faces_sharing_edge},
        \code{fix}/\code{refine}), all the \code{Gmake_*} constructors and the
        boolean operations. New test \code{tests/geo/test_occ_impl.py}: 39/39
        against real geometry with known expected values.
  \item \textbf{Phase 3 \textemdash{} Integration pipeline.}
        \code{tests/test_cadtocsg.py} passes 50/50 under
        \code{GEOUNED_CAD_ENGINE=occ}, identical to the FreeCAD baseline, after
        closing 6 real gaps (among them \code{Gload_step_labels} via XCAF, a
        periodicity bug in \code{get_join_cone_cyl} reachable under FreeCAD too,
        and the \code{GeomAPI_ProjectPointOnSurf} fallback).
  \item \textbf{Phase 4 \textemdash{} \code{Solidos/} corpus.} Differential scan
        of 98/99 files against a fresh FreeCAD baseline; 7 more gaps closed.
        \code{rev_pipe.stp} confirmed: FreeCAD fails the full decompose+convert
        path (\code{FuseEdges}), \code{occ} succeeds with
        \code{RoundC=2, MultiRoundC=1}.
\end{enumerate}

\subsection{Third engine: OCP (pybind11)}

Added alongside FreeCAD and pythonocc-core, not as a replacement. OCP is a
pybind11 binding of OCCT that is more direct and more actively maintained than
pythonocc-core's SWIG. Go/no-go validation against \code{rev_pipe.stp} (2
pieces, volume conserved), full port structurally copied from
\code{_occ_impl.py} and retargeted (with API differences confirmed live: \code{_s}
suffix on static methods, no \code{.DownCast()}, swapped argument order in
\code{ShapeUpgrade_UnifySameDomain}, \dots). \code{tests/geo/test_ocp_impl.py}
39/39 on the first run; \code{tests/test_cadtocsg.py} 50/50;
\code{tests/test_csgtocad.py} 2/2. Both pyOCC engines are developed in parallel:
a fix in one is ported to the other.

\paragraph{The performance gap, reported honestly.}
On the motivating case (\code{hylife-v06.stp} solid 17), FreeCAD took 2.3\,s and
pythonocc-core 19.2\,s. It was first attributed to the per-call cost of the SWIG
binding (real: pythonocc-core exposes the raw OCCT API without the C++
convenience layer FreeCAD has). OCP on the same case: 18.4\,s \textemdash{} only
$\sim$4\% faster, nowhere near the 2\textendash 6$\times$ of the point-conversion
micro-benchmarks. The real cause turned out to be \textbf{something else}:
\code{splitTolerance} with an implicit value of 0.0 made OCCT 7.9.3 compute and
then discard a degenerate 15-fragment split. With the default corrected to
$10^{-4}$ under \code{occ}/\code{ocp}, the same case drops to \textbf{1.2\,s}
\textemdash{} faster than FreeCAD.

% =============================================================================
\section{GEOReverse migration}

GEOReverse (CsgToCad) reconstructs CAD from CSG and historically depends on
FreeCAD's \code{.FCStd} document format, not just its geometry kernel. It was
initially permanently out of scope; on the current branch
(\code{georeverse-migration}) it has received the same treatment as the forward
pipeline.

\begin{itemize}
  \item \textbf{Decoupling.} The GEOReverse core (\code{Objects.py},
        \code{buildSolidCell.py}, \code{splitFunction.py}, \code{buildCAD.py},
        \code{MCNPinput.py}) imports \code{geo} directly, like the forward
        pipeline. A \code{matrix_utils.py} module holds the
        numpy\,$\leftrightarrow$\,native-matrix bridge.
        \code{BOPTools.SplitAPI.slice} $\rightarrow$ \code{Gsplit};
        \code{FuseSolid} (triplicated byte-for-byte) $\rightarrow$ a single
        \code{fuse_solids} with \code{GSolid.refine()}'s volume-invariance
        check.
  \item \textbf{Engine-independent \code{export_cad}.}
        \code{export_cad(output_filename="", format="stp")} accepts a
        \code{str} or a list of formats, validates against the engine's
        supported formats and dispatches by \code{CAD_ENGINE}. Real behaviour
        change: \code{.FCStd} is no longer written unconditionally, only when
        \code{"fcstd"} is explicitly requested.
  \item \textbf{Full FreeCAD/pyOCC symmetry.}
        \code{GEOReverse/Modules/__init__.py} dispatches between
        \code{_freecad_impl.py} (export via a FreeCAD document, \code{makeTree})
        and \code{_ocp_impl.py}/\code{_occ_impl.py} (real export via XCAF:
        \code{STEPCAFControl_Writer}, a
        Universe$\rightarrow$Material$\rightarrow$Cell hierarchy with names and
        nesting). The 6 exotic quadric surfaces were reorganised into a
        \code{geo_quadrics} package with the same dispatch pattern.
  \item \textbf{Per-material colour (\code{occ}/\code{ocp} engine).} Each solid
        in the exported STEP is coloured by its \code{MAT} value via XCAF: the
        same colour for the same material, a \code{tab10} palette up to 10
        materials and golden-angle rotation beyond that, with no collisions.
        Colour on the FreeCAD engine is out of scope (its colour API lives in
        \code{Gui.ViewProvider}, with no headless path).
  \item \textbf{IGES dropped.} It was verified live that
        \code{IGESControl_Writer} does not preserve solids through IGES in this
        OCCT build (a solid comes back as trimmed planar faces with no shell
        topology). IGES import and export were abandoned.
\end{itemize}

Verification: \code{tests/test_csgtocad.py} is engine-aware and passes 2/2 under
all three engines against real per-solid volume assertions ($\pm 10^{-6}$).
Under \code{occ}/\code{ocp} there remain 2 pre-existing failures (GEOReverse
produces 2/1 solids where 4/5 are expected in the \code{cylinder_box} STEP),
confirmed not introduced by this migration and not tied to the \code{Gsplit}
work.

% =============================================================================
\section{CAD defect detection and repair}
\label{sec:defects}

A new load-time layer (\code{Gload_and_process_step} $\rightarrow$
\code{Gcheck_and_repair}) that, for each solid: applies a \code{fix(1e-6)},
checks for known defects and, if any, tries to repair them with the most
appropriate method; only if the repair fails is the user-chosen mode consulted.
The \code{corrupted_solids} parameter has two modes, \code{"stop"} (default) and
\code{"remove"}; the old \code{invalid_solids} parameter was merged into it.

\begin{longtable}{@{}L{0.19\textwidth} L{0.29\textwidth} L{0.28\textwidth} L{0.15\textwidth}@{}}
\toprule
\textbf{Defect} & \textbf{Detection} & \textbf{Repair} & \textbf{Status} \\
\midrule
\endhead

\textbf{Split boundary ring} (duplicated micro-trim / collapsed step, spherical
variant) &
\code{find_split_ring_faces}: a sliver-scale analytic curved face that also
touches a pathologically short edge. \code{count_split_ring_pairs} counts pairs
of near-coincident, concentric circular edges. &
\code{Gcollapse_split_rings} (ocp/occ): removes the ``riser'' faces, re-sews the
shell, rebuilds the solid. Acceptance by 3 checks: validity +
$|\Delta V|/V < 3\times10^{-4}$ + the ring pairs must decrease. &
\statusok\newline \code{barrel bottom.stp} $\rightarrow$ tally 0.9997 \\
\addlinespace

\textbf{Cylindrical collapsed step} (beltline family) &
The same split-ring detector fires; \code{near_surface_pair} identifies two
cylinders of near-equal radius bridged by a sliver band. &
\code{Gsliver_heal} (ocp/occ), version 0: removes the sliver faces and the
smaller of the near pair, re-trims the freed quadric to its real radius, sews
\emph{last}. Gate: validity + $|\Delta V| < 5\times10^{-4}$. &
\statusok\newline \code{LR.stp} $\rightarrow$ tally 0.9989 (was 0.0 / 24 lost
particles) \\
\addlinespace

\textbf{Phantom cut} (two volumetric regions touching only along edges
\textemdash{} the migration's founding motivation) &
\code{_separate_edge_joined_components}: non-manifold edges (shared by $\neq$ 2
faces) $\rightarrow$ a face-adjacency graph over manifold edges only
$\rightarrow$ if there are $\geq$ 2 connected components and the non-manifold
edges were the load-bearing connection. &
Sews each component separately (no donor faces: a real phantom cut already
carries each region's closed boundary). Gate: $\geq$ 2 pieces, all valid, summed
volume conserved to $10^{-6}$. &
\statusdone\newline ocp/occ $\cdot$ \code{Piece_1}+\code{Tool_1} $\rightarrow$ 2
valid solids (was 1 invalid merged solid). FreeCAD cannot do it. \\
\addlinespace

\textbf{Malformed face wire} (\code{UnorientableShape}) &
\code{check_solid_defects} reports real invalid topology;
\code{BRepCheck_Analyzer} flags the face, no sliver edge. &
None. \code{ShapeFix_Shape}, \code{ShapeFix_Face}, wire partitioning and
rebuilding were all tried \textemdash{} they all fail or drift the volume $>$
1.7\%. &
\statusnone\newline \code{L4_support.stp} moved to
\code{Solidos/Detected_corrupted/}, excluded from conversion \\
\bottomrule
\end{longtable}

A related change: \code{check_solid_defects} no longer rejects a solid for a
sliver edge on a \emph{well-formed} face (normal \code{CharacteristicWidth})
\textemdash{} the decomposition pipeline already tolerates sliver edges on sound
faces, and no repair removes such an edge without a sliver face to anchor on.
The defect recipes are kept in a living catalogue
(\code{reference_cad_defect_recipes.md}): observable symptom, a general
detection signature, root cause, a repair if it is fast and reliable, and a
minimal STEP fixture with d1suned verification.

% =============================================================================
\section{Current status and open work}
\label{sec:pending}

\subsection{Consolidated status}

\begin{itemize}
  \item All three engines pass their suites: \code{freecad} 158/158, \code{occ}
        128/128, \code{ocp} 128/128 (+2 pre-existing GEOReverse failures under
        the pyOCC engines).
  \item Corpus \code{Solidos/test_models}: 93.6\% $<$ 2$\sigma$, 0\% $>$
        3$\sigma$, 0 lost particles across 143 d1suned runs. \code{occ} and
        \code{ocp} are byte-for-byte equivalent over the whole corpus.
  \item \code{ocp} is the default engine when \code{GEOUNED_CAD_ENGINE} is
        unset.
\end{itemize}

\subsection{Open}

\begin{longtable}{@{}L{0.62\textwidth} L{0.32\textwidth}@{}}
\toprule
\textbf{Item} & \textbf{Nature} \\
\midrule
\endhead
Volume discrepancy in the GEOReverse round-trip on \code{hylife-v06.stp} cell 1
(reconstructed CAD 1.401$\times$ vs d1suned tally 1.119$\times$). & GEOReverse,
not traced to the bottom. \\
\code{test_csgtocad.py} 2 failures under \code{occ}/\code{ocp} (2/1 solids where
4/5 are expected). & GEOReverse, pre-existing, narrowed to
\code{Objects.py::CellObj.buildShape}. \\
\code{AdjacentMultiplanePlanes} needs extending from RevCC to MultiRoundCorner. &
Coverage improvement, open since 2026-08-21. \\
The 6 exotic quadrics (\code{Gmake_elliptic_cone}, \code{Gmake_hyperboloid},
\dots) are stubs in GEOReverse's pyOCC backends. & Port pending; no fixture
exercises them yet. \\
\code{Big_complex_cell/modelcell_cut1.stp} and \code{modelCell_670000.stp} lose
particles. & Not traced; excluded from corpus scans by convention. \\
\code{hylife-v06.stp} full 372-solid round-trip: slow decomposition in
\code{meta_list[45]} (duplicated face fragments, $O(n^2)$ cost in
\code{same_faces}). & Root cause identified, not optimised, deprioritised by the
user. \\
\code{ConeSphere.stp}: native crash (access violation) under
\code{occ}/\code{ocp} in \code{ShapeUpgrade_UnifySameDomain}; translates clean
under \code{freecad}. & \statuslimit{} No OCCT knob avoids it; documented in the
docstrings. \\
Reorganisation and dedup of the \code{Solidos/} fixture tree. & Partial
(\code{test_models/} + a few \code{duplicates_removed/}). \\
\bottomrule
\end{longtable}

Explicit user priority: finish cleaning up the known bugs in the forward
pipeline (GEOUNED) before resuming the GEOReverse items.

% =============================================================================
\section{Appendix: environment and tooling}

\begin{center}
\small
\begin{tabular}{@{}L{0.20\textwidth} L{0.24\textwidth} L{0.18\textwidth} L{0.30\textwidth}@{}}
\toprule
\textbf{conda env} & \textbf{Binding} & \textbf{OCCT $\cdot$ Python} & \textbf{Use} \\
\midrule
(default python) & FreeCAD 1.1.1 & 7.8.1 $\cdot$ 3.11 & \code{GEOUNED_CAD_ENGINE=freecad} \\
\code{pyoccenv} & pythonocc-core (SWIG) & 7.9.3 $\cdot$ 3.12 & \code{GEOUNED_CAD_ENGINE=occ} \\
\code{ocpenv} & OCP (pybind11) & 7.9.3 $\cdot$ 3.12 & \code{GEOUNED_CAD_ENGINE=ocp} (default) \\
\bottomrule
\end{tabular}
\end{center}

\begin{itemize}
  \item \textbf{FreeCAD and pyOCC cannot coexist in one process}
        (\code{Module use of python311.dll conflicts}): FreeCAD ships Python
        3.11, the pyOCC environments use 3.12. Any pyOCC-side test uses that
        environment's \code{python.exe} and never mixes in the same invocation.
  \item The pyOCC environments needed the BLAS/LAPACK provider switched from MKL
        to OpenBLAS
        (\code{conda install -c conda-forge --override-channels "libblas=*=*openblas" ...}):
        MKL's LAPACK conflicts with \code{occt}/\code{tbb} in the same process
        and causes an access violation in \code{numpy.linalg}.
  \item The MCNP volume check runs on \textbf{d1suned} (an in-house MCNP
        variant) through WSL. The full pipeline, scripts and results are
        archived under WSL (Ubuntu-22.04) at
        \code{~/work/taller/SolidTestMCNP/scripts/} (parallel conversion and
        d1suned per engine: \code{run_all_conversions_*_parallel.py},
        \code{run_test_models_*_d1suned.sh}, \code{analyze_test_models_*.py}).
  \item Working-tree convention: \code{Solidos/} holds STP only;
        \code{Solidos/test_models/} is the curated regression set;
        \code{../calculos_CLAUDE/} holds every calculation run;
        \code{Solidos/BadCAD_decomposition/},
        \code{Solidos/Detected_corrupted/}, \code{Solidos/no_convierte/},
        \code{Solidos/convierte_bad_volume/} are triage folders.
\end{itemize}

\vfill
\begin{center}\footnotesize\itshape
Report synthesised from the GEOUNED project working log \code{CLAUDE.md}. Branch
\code{georeverse-migration}. Date: 2026-09-06.\\
Project code-style preference: communication is in Spanish; all code, comments,
docstrings and identifier names are in English.
\end{center}

\end{document}
